% УВОД, без номер. Цел, задачи, обхват, състояние.
\section*{УВОД}
\label{sec:introduction}
\addcontentsline{toc}{section}{УВОД}

Публичните облачни платформи се представят като взаимозаменяеми: една и съща услуга
може да бъде разгърната при кой да е доставчик, а изборът се свежда до цена. Проверката
на това твърдение изисква едно и също натоварване да бъде разгърнато и измерено на две
места при съпоставими условия, което рядко се прави.

Проектът изгражда \term{Stratus}: изчислително интензивна услуга в контейнер, обратен
посредник, генератор на синтетично натоварване и контролер за мащабиране, всичките на
C++20, заедно с описание на инфраструктурата, което насочва същата услуга към Amazon Web
Services или Microsoft Azure чрез стойността на една променлива.

\textbf{Целта} е да се провери доколко едно декларативно описание на инфраструктурата е
достатъчно за разгръщане на една и съща услуга при двама независими доставчика, и да се
измери разликата между тях по закъснение, поведение при мащабиране и цена за единица
работа.

\textbf{Задачите} са четири: кратка теоретична рамка и обосновка на направените
технологични избори; проектиране, при което платформено зависимата част е сведена до
преброими места; програмна реализация на услугата, наблюдението и мащабирането;
методика за измерване и провеждане на измерванията.

\textbf{Обхватът} е ограничен до един тип изчислително натоварване и до двама доставчици.
Извън обхвата остават хибридните и частните облаци, миграцията на състояние по време на
работа и управляваните бази от данни.

\textbf{Докъде е стигнала работата.} Този абзац стои в увода, защото определя как да се
четат всички числа по-нататък. Изградени и измерени са услугата, слоят за наблюдение,
контролерът за мащабиране и генераторът на натоварване, като измерванията са направени
върху локален стек от контейнери на една машина. Описанието на инфраструктурата за двамата
доставчици е написано и форматирано, но \textbf{не е прилагано}: няма данни за достъп до
сметка при нито един от тях и \code{terraform} не е инсталиран на използваната машина.
Затова в документа няма \emph{нито едно} измерване при AWS или Azure. Липсващите места са
означени явно, вместо да бъдат запълнени с правдоподобни стойности.
